\lecture{8.3}{Longitudinal G-Computation Formula}
\label{lec:8-3}

\subsection{Two-Timepoint Derivation}

Consider the observed data
\[
  O=(L_0,A_0,L_1,A_1,Y)
\]
and the target parameter
\[
  \psi(1,1)=\E\{Y(1,1)\}.
\]
Assume the three identifying conditions hold: consistency, sequential
exchangeability, and positivity. The formula repeatedly replaces counterfactual
conditional means with observed conditional means, then averages over the
relevant covariate distributions.

Start with the innermost observed outcome regression:
\[
  \E(Y\given A_1=1,L_1,A_0=1,L_0).
\]
By consistency, among individuals with $A_0=1$ and $A_1=1$, the observed outcome
equals $Y(1,1)$. Therefore
\[
  \E(Y\given A_1=1,L_1,A_0=1,L_0)
  =
  \E\{Y(1,1)\given A_1=1,L_1,A_0=1,L_0\}.
\]
By sequential exchangeability at time 1,
\[
  Y(1,1)\ci A_1\given L_0,A_0=1,L_1,
\]
so
\[
  \E\{Y(1,1)\given A_1=1,L_1,A_0=1,L_0\}
  =
  \E\{Y(1,1)\given L_1,A_0=1,L_0\}.
\]
Define
\[
  \Qbar_2(L_0,L_1)
  =
  \E(Y\given A_1=1,L_1,A_0=1,L_0).
\]
Then
\[
  \Qbar_2(L_0,L_1)
  =
  \E\{Y(1,1)\given L_1,A_0=1,L_0\}.
\]

Now average over the distribution of $L_1$ among those with $A_0=1$ and $L_0$:
\[
  \E\{\Qbar_2(L_0,L_1)\given A_0=1,L_0\}
  =
  \E\{Y(1,1)\given A_0=1,L_0\}.
\]
By sequential exchangeability at time 0,
\[
  Y(1,1)\ci A_0\given L_0,
\]
so
\[
  \E\{Y(1,1)\given A_0=1,L_0\}
  =
  \E\{Y(1,1)\given L_0\}.
\]
Define
\[
  \Qbar_1(L_0)=\E\{\Qbar_2(L_0,L_1)\given A_0=1,L_0\}.
\]
Finally, average over the baseline covariate distribution:
\[
  \E\{\Qbar_1(L_0)\}
  =
  \E\{Y(1,1)\}.
\]

\begin{identificationbox}
For two treatment times,
\[
  \E\{Y(1,1)\}
  =
  \E\left[
    \E\left\{
      \E(Y\given A_1=1,L_1,A_0=1,L_0)
      \given A_0=1,L_0
    \right\}
  \right].
\]
The outer expectation is over the marginal distribution of $L_0$.
\end{identificationbox}

\subsection{Iterated Conditional Expectations}

The formula is a direct application of the tower rule, but the conditioning sets
are chosen according to the intervention. The innermost regression describes the
outcome among individuals who followed the regime through the final treatment.
The next expectation averages over the covariate that would be observed after
following the regime through the previous time. Repeating this process moves
backward through time.

\begin{figure}[H]
\centering
\begin{tikzpicture}[
  node distance=.8cm,
  box/.style={draw, rounded corners=2pt, minimum width=3.2cm, minimum height=.72cm},
  arr/.style={-{Latex[length=2.2mm]}, thick}
]
\node[box] (q2) {$\Qbar_2(L_0,L_1)=\E(Y\given A_0=1,A_1=1,L_0,L_1)$};
\node[box, below=of q2] (q1) {$\Qbar_1(L_0)=\E\{\Qbar_2(L_0,L_1)\given A_0=1,L_0\}$};
\node[box, below=of q1] (q0) {$\psi(1,1)=\E\{\Qbar_1(L_0)\}$};
\draw[arr] (q2) -- node[right] {average over $L_1$} (q1);
\draw[arr] (q1) -- node[right] {average over $L_0$} (q0);
\end{tikzpicture}
\caption{Longitudinal g-computation as backward nested averaging.}
\label{fig:ch8-gcomp-nesting}
\end{figure}

\subsection{General $K$-Timepoint Formula}

For a fixed regime $\abar=(a_0,\ldots,a_K)$, define
\[
  \Qbar_{K+1}(\Lbar_K)
  =
  \E(Y\given \Abar_K=\abar_K,\Lbar_K).
\]
Then recursively define, for $k=K,K-1,\ldots,1$,
\[
  \Qbar_k(\Lbar_{k-1})
  =
  \E\{\Qbar_{k+1}(\Lbar_k)
    \given \Abar_{k-1}=\abar_{k-1},\Lbar_{k-1}\}.
\]
The target mean is
\[
  \psi(\abar)=\E\{\Qbar_1(L_0)\}.
\]

Equivalently, the formula can be displayed as nested expectations:
\[
  \E\{Y(\abar)\}
  =
  \E\left[
  \E\left[
  \cdots
  \E\left[
  \E(Y\given \Abar_K=\abar_K,\Lbar_K)
  \given \Abar_{K-1}=\abar_{K-1},\Lbar_{K-1}
  \right]
  \cdots
  \given A_0=a_0,L_0
  \right]\right].
\]

\subsection{Sequential Outcome Regressions}

The functions $\Qbar_{K+1},\Qbar_K,\ldots,\Qbar_1$ are called sequential outcome
regressions. They are not all regressions of $Y$. Only $\Qbar_{K+1}$ is the
conditional mean of the observed outcome. The earlier regressions use later
predicted values as pseudo-outcomes.

For $K=1$, the sequence is:
\[
  \Qbar_2(L_0,L_1)=\E(Y\given A_0=a_0,A_1=a_1,L_0,L_1),
\]
\[
  \Qbar_1(L_0)=\E\{\Qbar_2(L_0,L_1)\given A_0=a_0,L_0\},
\]
\[
  \psi(a_0,a_1)=\E\{\Qbar_1(L_0)\}.
\]

\subsection{Interpretation of Nested Expectations}

Each expectation performs one of two tasks. First, it replaces the observed
treatment assignment with the regime value. Second, it averages over covariates
that would naturally arise after earlier parts of the regime.

For the regime $(1,1)$, the inner regression asks: among people who received
treatment at both times and had a given covariate history, what was the mean
outcome? The next regression asks: among people who received treatment at time
0 and had a given $L_0$, what covariate histories $L_1$ occurred, and what
outcomes would those histories imply if treatment also occurred at time 1? The
final average transports these predicted means to the baseline target population.

\begin{keyidea}
Longitudinal g-computation is standardization over time. It standardizes the
final outcome regression over the sequence of covariate distributions generated
under the intervention history.
\end{keyidea}
